% ---
% titre: Échelle (carte, plan, modèle réduit)
% theme: FA
% chapitre: Proportionnalité
% sous_chapitre: Échelle
% niveau: [10e]
% type: EVAL
% tags: [c-proportionnalite,c-echelle,p-question-TS]
% difficulte: 1
% corrige: true
% ---

\question
Qu'elle est l'échelle?

\begin{parts}
\vspace{20pt}\part[3] Sur une carte nationale à l'échelle 1:\,25\,000, la distance entre deux refuges mesure 7,2\,cm. Quelle est la distance réelle entre ces deux refuges ? Donne ta réponse en kilomètres.

\vspace{10pt}{\footnotesize\raggedleft Espace pour ta démarche et tes calculs\par}
\begin{solutionorgrid}[200pt]

\vspace{5pt}
\def\arraystretch{1.5}
\begin{tabularx}{0.6\textwidth}{l|C|C|}
Carte [cm]&1&7,2\\
\hline
Réalité [cm]&25\,000&\\
\end{tabularx}
\def\arraystretch{1}

\hfill\textbf{(0,5\,pt)} pour la construction du tableau

\vspace{15pt}$7,2 \cdot 25\,000 : 1 = 180\,000$\,cm \hfill calcul \textbf{(0,5\,pt)} et réponse \textbf{(0,5\,pt)}

\vspace{5pt}$180\,000\,\text{cm} = 1\,800\,\text{m} = 1,8\,\text{km}$ \hfill \textbf{(1\,pt)}
\end{solutionorgrid}

\begin{tcolorbox}[enhanced jigsaw, interior hidden, colframe=box, parbox=false]%
Réponse : 
\sol{\vspace{20pt}}{La distance réelle entre les deux refuges est de 1,8\,km. \textbf{(0,5\,pt)}}
\end{tcolorbox}

\vspace{20pt}\part[2] Sur le plan d'un appartement, la longueur du salon mesure 9\,cm. Dans la réalité, elle mesure 4,5\,m. Quelle est l'échelle de ce plan ?

\vspace{10pt}{\footnotesize\raggedleft Espace pour ta démarche et tes calculs\par}
\begin{solutionorgrid}[170pt]

\vspace{5pt}$4,5\,\text{m} = 450\,\text{cm}$ \hfill \textbf{(0,5\,pt)}

\vspace{5pt}
\def\arraystretch{1.5}
\begin{tabularx}{0.6\textwidth}{l|C|C|}
Plan [cm]&9&1\\
\hline
Réalité [cm]&450&\\
\end{tabularx}
\def\arraystretch{1}

\hfill\textbf{(0,5\,pt)} pour la construction du tableau

\vspace{15pt}$450 \cdot 1 : 9 = 50$ \hfill \textbf{(0,5\,pt)}
\end{solutionorgrid}

\begin{tcolorbox}[enhanced jigsaw, interior hidden, colframe=box, parbox=false]%
Réponse : 
\sol{\vspace{20pt}}{L'échelle du plan est 1:\,50. \textbf{(0,5\,pt)}}
\end{tcolorbox}

\vspace{20pt}\part[2] Un modèle réduit de voiture est à l'échelle 1:\,18. La voiture réelle mesure 432\,cm de long. Quelle est la longueur du modèle réduit ?

\vspace{10pt}{\footnotesize\raggedleft Espace pour ta démarche et tes calculs\par}
\begin{solutionorgrid}[170pt]

\vspace{5pt}
\def\arraystretch{1.5}
\begin{tabularx}{0.6\textwidth}{l|C|C|}
Modèle [cm]&1&\\
\hline
Réalité [cm]&18&432\\
\end{tabularx}
\def\arraystretch{1}

\hfill\textbf{(0,5\,pt)} pour la construction du tableau

\vspace{15pt}$432 \cdot 1 : 18 = 24$\,cm \hfill calcul \textbf{(0,5\,pt)} et réponse \textbf{(0,5\,pt)}
\end{solutionorgrid}

\begin{tcolorbox}[enhanced jigsaw, interior hidden, colframe=box, parbox=false]%
Réponse : 
\sol{\vspace{20pt}}{Le modèle réduit mesure 24\,cm de long. \textbf{(0,5\,pt)}}
\end{tcolorbox}
\end{parts}
